\section{Related Work}\label{sec:related}

\paragraph{Text-to-motion generation.} Human text-to-motion has moved from VAEs on
HumanML3D~\citep{guo2022generating} to diffusion in pose space~\citep{tevet2023mdm,zhang2024motiondiffuse}
or a latent~\citep{chen2023mld}, distilled for real time~\citep{dai2024motionlcm}, and to token
models over a VQ-VAE~\citep{zhang2023t2mgpt,guo2024momask}, scaled by larger
corpora~\citep{lin2023motionx,fan2025motionmillion} and skeleton-aware latent
attention~\citep{hong2025salad}. All fix one kinematic tree; \method{} keeps the pose-space
diffusion and text conditioning and removes the fixed skeleton.

\paragraph{Motion across skeleton topologies.} Cross-skeleton motion has mostly meant
\emph{retargeting} between two known rigs~\citep{villegas2018nkn,aberman2020skeleton} or
skeleton-agnostic embeddings such as SAME~\citep{lee2023same} and SATA~\citep{zhang2026sata}, whose
text-conditioned generation is a downstream demonstration on HumanML3D. For generation on arbitrary
topologies, SinMDM~\citep{raab2024sinmdm} fits one motion at a time; AnyTop~\citep{gat2025anytop}
generates, without text, across the skeletons of the Truebones zoo~\citep{truebones}; UniMoGen~\citep{khani2025unimogen} is skeleton-agnostic and controlled by
style and trajectory. Text-conditioned animal motion has been approached from a fixed 35-joint
SMAL template by OmniMotionGPT~\citep{yang2024omnimotiongpt}, from a shared human--animal topology
by X-MoGen~\citep{wang2025xmogen} (115 species), across topologies by OmniZoo~\citep{chen2025omnizoo}
(a topology-aware skeleton embedding and 140 species), on a unified 30-joint skeleton by AniMo~\citep{wang2025animo}, whose AniMo4D
captions our corpus carries, on user-defined human skeletons by OmniSkel~\citep{liu2025omniskel},
on re-captioned Truebones data by G-MDM~\citep{lee2025dragon},
which generates Euler joint rotations without root translation, and, concurrently with this work, by
SAMoR~\citep{samor2026}, which pools any skeleton into eight part tokens over joint positions and
velocities and recovers rotations by inverse kinematics, and by UniMate~\citep{mou2026unimate}. \method{} is text-conditioned; it
denoises a per-joint representation any rig is converted to without a template, retargeting or inverse
kinematics; it embeds natural-language \emph{descriptions} of the joints with LLM2Vec; and it normalises
every (rig, joint, channel) cell with its own statistics, so that
hundreds of rigs of very different proportions share one network.

\paragraph{Objective, encoders, adaptation and evaluation.} \method{} is a DiT-style
denoiser~\citep{peebles2023dit} trained with flow matching~\citep{lipman2023flow,liu2023rectified},
classifier-free guidance on the text stream~\citep{ho2022cfg} and query--key
normalisation~\citep{dehghani2023vit22b}; its grouped per-channel-family objective follows
Kimodo~\citep{rempe2026kimodo}. Captions and joint descriptions are embedded with
LLM2Vec~\citep{behnamghader2024llm2vec} rather than CLIP~\citep{radford2021clip} or
T5~\citep{raffel2020t5}. New rigs are handled zero-shot through the skeleton prompt, in the spirit of
in-context conditioning for image diffusion~\citep{wang2023promptdiffusion}, and with a per-rig
LoRA~\citep{hu2022lora} when they lie far from the library. Retrieval metrics follow the contrastive text--motion evaluators
of~\citet{guo2022generating} and TMR~\citep{petrovich2023tmr}, which are tied to one skeleton; ours
is skeleton-aware (\S\ref{sec:method:eval}).
